\chapter{Comparative Analysis}
\label{Comparison:chap:chap8}
\graphicspath{{17.Chapter8/Images/}}
\raggedbottom

This chapter places FIP-FL alongside the Byzantine-resilient FL defences that were already covered in Chapter~\ref{Literature:chap:chap3}, with the auditor-based baseline of Yazdinejad et al.~\cite{yazdinejad2024robust} taking centre stage. The comparison runs over six axes, privacy leakage, computation, communication, robustness under non-IID skew, engineering cost, and the canonical MNIST non-IID targeted benchmark.

\section{Privacy}
\label{sec:comp-privacy}

The auditor-based scheme has to decrypt cluster-level statistics every round: a projected Gaussian-mixture model along with the Mahalanobis-distance statistics that the filter uses. Even though no individual gradient is ever reconstructed, the cluster means and covariances themselves still leak structural information about the client population. In a non-IID deployment, those statistics can correlate with demographics, hospital specialisation, keyboard language, or any number of other sensitive client traits.

FIP-FL is much more spare. The only thing decrypted per client per round is the scalar FIP score $s_i^{(t)}=\langle \mathbf{g}_i^{(t)},\mathbf{R}^{(t)}\rangle$. There are no cluster means, no covariance matrices, no projection coefficients, no pairwise distances. The privacy improvement is categorical rather than incremental: the verifier sees a single alignment number, and the model update is still derived only from an aggregate.

\section{Computational Complexity}
\label{sec:comp-compute}

\begin{table}[H]
\centering
\caption{Per-round computational cost. $N$: clients; $|\mathcal{A}|$: accepted clients; $d$: model dimension; $B$: batch size; $E$: local epochs.}
\label{tab:comp-computation}
\begin{tabular}{@{}lll@{}}
\toprule
\textbf{Entity} & \textbf{Auditor-based~\cite{yazdinejad2024robust}} & \textbf{FIP-FL} \\
\midrule
Client              & $\mathcal{O}(dBE)$ & $\mathcal{O}(dBE)$ \\
Verifier / Auditor  & $\mathcal{O}(NKD+Nd\log d)+\mathcal{O}(d^3)$ & $\mathcal{O}(Nd+N\log N)$ \\
Aggregation Server  & $\mathcal{O}(dN\log N)$ & $\mathcal{O}(|\mathcal{A}|d)$ \\
\bottomrule
\end{tabular}
\end{table}

\subsection{Time-Complexity Interpretation}
\label{sec:comp-time-interpretation}

The computational comparison in Table~\ref{tab:comp-computation} shows that the main saving of FIP-FL appears at the verifier stage. The client-side cost remains unchanged because both schemes require each client to perform ordinary local SGD on its own data. The difference lies in how the submitted updates are verified and aggregated after local training.

\noindent\textbf{Client-side training.}
The auditor-based PPFL scheme and FIP-FL both have client-side complexity $\mathcal{O}(dBE)$. Both schemes have the same local learning cost because each client trains on mini-batches for $E$ local epochs.

\noindent\textbf{Verification basis.}
The auditor-based PPFL scheme verifies updates through distributional modelling of submitted gradient statistics, while FIP-FL verifies updates through scalar alignment with the bootstrap reference direction. FIP-FL reduces verification from high-dimensional distributional analysis to one scalar score per client.

\noindent\textbf{Verifier / Auditor operations.}
The auditor-based scheme uses projection, Gaussian-mixture modelling, Mahalanobis-distance scoring, and covariance-related computation. FIP-FL uses homomorphic functional inner product computation and bootstrap score filtering. The auditor-based scheme performs richer statistical modelling, whereas FIP-FL performs one encrypted inner product per client.

\noindent\textbf{Verifier / Auditor complexity.}
The auditor-based verifier has complexity $\mathcal{O}(NKD+Nd\log d)+\mathcal{O}(d^3)$, while the FIP-FL verifier has complexity $\mathcal{O}(Nd+N\log N)$. Under the asymptotic model used in this thesis, the additional $\mathcal{O}(d^3)$ term accounts for covariance-inversion style computation in the auditor-side distributional model.

\noindent\textbf{Aggregation cost.}
The auditor-based aggregation cost is $\mathcal{O}(dN\log N)$, while the FIP-FL aggregation cost is $\mathcal{O}(|\mathcal{A}^{(t)}|d)$. FIP-FL aggregates only the accepted ciphertexts after verification, while rejected ciphertexts are dropped.

\noindent\textbf{Scalability implication.}
In the auditor-based PPFL scheme, cost increases with distributional modelling and covariance-related computation. In FIP-FL, verification is linear in the model dimension at the verifier stage. FIP-FL is therefore more scalable for high-dimensional updates because verification avoids cluster-level statistics and covariance estimation.

Here, $N$ is the number of clients, $d$ is the model dimension, $B$ is the batch size, $E$ is the number of local epochs, $K$ is the number of mixture components, $D$ is the auditor-side feature or projection dimension, and $\mathcal{A}^{(t)}$ is the accepted client set in round $t$.

Overall, this comparison shows that FIP-FL keeps the same client-side learning cost as the auditor-based baseline but replaces auditor-side distributional modelling with a scalar encrypted inner-product audit. For each client, the verifier computes one homomorphic functional inner product with the fixed bootstrap reference direction, giving $\mathcal{O}(Nd)$ score-computation cost. The bootstrap audit then filters scalar scores, contributing $\mathcal{O}(N\log N)$. After verification, only accepted ciphertexts are aggregated, giving $\mathcal{O}(|\mathcal{A}^{(t)}|d)$. Thus, FIP-FL shifts verification from high-dimensional statistical auditing to linear-time encrypted alignment testing, making the verifier stage easier to scale and parallelise across clients.

\section{Communication Complexity}
\label{sec:comp-comm}

\begin{table}[H]
\centering
\caption{Per-round communication cost by directed link.}
\label{tab:comp-comm}
\begin{tabular}{@{}lll@{}}
\toprule
\textbf{Link} & \textbf{Auditor-based~\cite{yazdinejad2024robust}} & \textbf{FIP-FL} \\
\midrule
Client $\to$ Verifier & $\mathcal{O}(Nd)$ & $\mathcal{O}(Nd)$ \\
Verifier $\to$ Server & $\mathcal{O}(L(Z+N))$ & $\mathcal{O}(|\mathcal{A}|d)$ \\
Server $\to$ Clients  & $\mathcal{O}(d)$ & $\mathcal{O}(d)$ \\
\bottomrule
\end{tabular}
\end{table}

The first leg of communication, the client-to-verifier payload, is the same in both designs, since both schemes need the encrypted updates to reach the verifier somehow. The savings show up after the audit. FIP-FL never forwards a rejected ciphertext, so the verifier-to-server traffic shrinks with the rejection rate, and the extra auditor-to-server hop that the auditor-based pipeline relies on simply does not exist here.

\section{Robustness and Engineering Trade-Offs}
\label{sec:comp-robust-engineering}

The detector inside the auditor-based scheme leans on the assumption that the honest-client cluster is both tighter and more populous than the malicious cluster, an assumption that gets noticeably weaker once label skew is in the mix. FIP-FL takes a different tack: it tests, for each individual update, whether the update is aligned with the trusted reference direction and whether its score passes the bootstrap audit. The whole admissibility check stays one-dimensional, and that one-dimensional behaviour is what carries the audit through the heterogeneous cases.

A useful side-effect of this design is that it limits cascade behaviour after a hard round. Only the updates that clear the bootstrap audit ever enter the aggregate, and the bootstrap reference remains fixed across rounds, which prevents a single noisy round from completely re-orienting the audit. The Phase~II FIP computations are independent across clients, so they parallelise cleanly, on the 25-core workstation described in Section~\ref{sec:setup-impl} this gave a wall-clock speed-up of about $13\times$. The single-reference design also avoids the extra homomorphic inner products that a multi-reference cone would have demanded.

\section{Scheme-versus-Scheme Accuracy Comparison}
\label{sec:comp-accuracy}

\begin{table}[H]
\centering
\caption{Accuracy on MNIST non-IID under a $50\%$ malicious-client rate. Entries are percentages; baseline values are from~\cite{yazdinejad2024robust}.}
\label{tab:comp-accuracy}
\begin{tabular}{@{}lcc@{}}
\toprule
\textbf{Scheme} & \shortstack{\textbf{Targeted}\\\textbf{target-class accuracy}\\\textbf{($T_{\text{acc}}$)}} & \shortstack{\textbf{Untargeted}\\\textbf{overall accuracy}\\\textbf{(ACC)}} \\
\midrule
Krum~\cite{blanchard2017machine}              & 71.4 & 78.3 \\
Auror~\cite{shen2016auror}                    & 32.5 & 49.9 \\
PEFL~\cite{liu2021privacy}                    & 46.4 & 57.4 \\
Sybils / FoolsGold~\cite{fung2020limitations} & 89.7 & 89.5 \\
FL-Trust~\cite{cao2021fltrust}                & 37.8 & 43.4 \\
ShieldFL~\cite{ma2022shieldfl}                & 90.1 & 89.4 \\
Auditor-based~\cite{yazdinejad2024robust}     & 96.5 & 96.3 \\
\midrule
\textbf{FIP-FL (this thesis)}                 & \textbf{97.69} & 88.82 \\
\bottomrule
\end{tabular}
\end{table}

In the table above, FIP-FL records the best target-class accuracy ($T_{\text{acc}}$) of every targeted entry, and it does so while resting on weaker trust assumptions than the auditor-based row. Its untargeted overall accuracy (ACC) is competitive with FoolsGold and ShieldFL but does not quite reach the auditor-based result, and that gap is the clearest empirical place where future work could push further. For the purposes of this thesis the relevant point is the qualitative one: FIP-FL produces the strongest target-class recovery without decrypting any cluster statistics and without bringing in an external auditor.
